% ============================================================
% Chapter 5: Development
%   Section: News Acquisition: Discovery and Extraction
%     Subsection: What Discovery Keeps
% Included from chapters/05-development/02-news-acquisition/news-acquisition.tex
% ============================================================

\subsection{What Discovery Keeps}
\label{subsec:scraping-filters}

A feed or a section page lists far more than articles worth offering:
section indexes, live blogs, opinion columns, sport, items from years
ago and articles already analysed. Every link is filtered before any
article page is fetched, which is the cheapest point at which to say
no. Table~\ref{tab:discovery-filters} summarises the rules; each was
added after a concrete case, given in the table.

\begin{table}[ht]
\centering
\small
\caption{Rules applied to discovered links, and the case behind each.}
\label{tab:discovery-filters}
\begin{tabularx}{\linewidth}{@{}p{3.3cm}X@{}}
\toprule
Rule & What is dropped, and why \\
\midrule
Article shape
& A link must look like an article: a date in its path
  (\texttt{/2026/09/23/}, \texttt{/20260923/}), a slug of at least four
  words, or an \texttt{/article/} or \texttt{/story/} segment. The
  previous English section-name patterns had kept 0 of El Mundo's 26
  feed items, 5 of El País's 149 and 0 of NASA's 10. \\
Off-mission sections
& Sport, celebrity, horoscopes, lotteries, betting, weather, crime
  (\texttt{/sucesos/}) and El País's lifestyle magazine
  (\texttt{/icon/}), as whole path segments, so that
  \texttt{/sports-health/} (fitness) stays. Trigger: a match report
  scored above an analysis of Chinese industry leaving fossil fuels on
  the positive-impact score (0.34 against 0.28). \\
Live coverage and opinion
& Live pages are streams of updates, not articles; a column's claims
  are interpretation (\texttt{/opinion/}, \texttt{/commentisfree/}). \\
Index pages
& \texttt{/categoria/}, \texttt{/tag/}, \texttt{/especiales/},
  \texttt{/autor/}\ldots, whose slugs read like headlines. \\
Section navigation
& On a section page every link contains the section, so a link there
  must look like an article by its own shape: BBC's section pages had
  given \texttt{/culture/music} as a candidate. \\
Outside the section
& A section page's links outside its own section are dropped when it
  has any inside it: ABC's \texttt{/cultura/} page linked the day's
  election stories. \\
Another topic's section
& A link filed under a section of a topic not asked for, unless the
  item mentions a topic that was: elDiario files wind-power records
  under \texttt{/economia/}, and those stay. \\
Topic keywords
& English items must match a requested topic's section name or
  keywords in their title, summary or categories. Spanish items, only
  when the round is scoped, must contain a Spanish topic word in their
  title or categories (Spanish summaries of general feeds mention
  renewables or the political ``clima'' in campaign speeches). \\
Age
& Older than 30 days, by the feed's time or, without one, by a date in
  the link (including ABC's slug timestamps). Undated links stay. \\
Already stored
& The URL's host and path are already in the data lake, so tracking
  parameters do not make an old article new. \\
\bottomrule
\end{tabularx}
\end{table}

\paragraph{Keywords and language.} The topic keywords are English.
Matched against a Spanish feed they kept 5 of El País's 149 items at
random, so Spanish feeds were at first not filtered at all, and an
Environment round then listed El País's and elDiario's election
coverage. A Spanish keyword list is now read when, and only when, a
round is narrowed to some topics, and only against an item's title and
categories. Measured on the live feeds on 5~October~2026, Environment
kept 18 of El País's 159 items and 6 of elDiario's 109, nearly all on
topic. When every topic is asked for (the one-shot ingestion endpoint
and the labelling batch), Spanish feeds are not filtered, and topic is
left to the admission filter, which rejects before any language-model
call.

\paragraph{The age limit.} Thirty days was chosen from a sample of
6~October~2026: 38 of 363 dated candidates were older (24 of them from
the World Health Organization, back to 2025), while a 14-day limit
would also have cut nine MIT News research stories. A link with no
feed time can still be dated by its path, which is how ABC's mortgage
and income-tax calculators, stamped in May, stopped being offered as
news in October. That date is used only for the age check: a day is
not a publication time. The labelling batch has its own discovery and
is not cut by this limit; it prefers articles from the last 60 days
instead.

\paragraph{How many per source.} At most \texttt{perSource} new
articles are taken from each source per run (three by default); the
rest are reported as deferred rather than silently lost. Before that
cap, the mission screen reads up to five times as many of each
source's new articles and leaves out those nearer politics, crime,
disasters, war, celebrity, sport, markets or service pages than any
topic, so that a source's slots go to what passes. The screen, the
ranking of the candidates and the choice among them are part of
selection, described in Section~\ref{sec:design-selection}.
